
\documentclass{article}
\usepackage{colortbl}
\usepackage{makecell}
\usepackage{multirow}
\usepackage{supertabular}

\begin{document}

\newcounter{utterance}

\centering \large Interaction Transcript for game `eqbench', experiment `eq\_bench', episode 121 with reflexion{-}Qwen3.5{-}4B.
\vspace{24pt}

{ \footnotesize  \setcounter{utterance}{1}
\setlength{\tabcolsep}{0pt}
\begin{supertabular}{c@{$\;$}|p{.15\linewidth}@{}p{.15\linewidth}p{.15\linewidth}p{.15\linewidth}p{.15\linewidth}p{.15\linewidth}}
    \# & \multicolumn{2}{c}{Player} && \multicolumn{2}{c}{Game Master} \\
    \hline

    \theutterance \stepcounter{utterance}  
    & & & \multicolumn{4}{p{0.6\linewidth}}{
        \cellcolor[rgb]{0.9,0.9,0.9}{
            \makecell[{{p{\linewidth}}}]{
                \texttt{\tiny{[P1$\langle$GM]}}
                \texttt{You have attempted this same task before and did not succeed. These are the plans you wrote after those attempts.} \\
\texttt{Trial \#1: New plan: I misinterpreted the emotional dynamic of the dialogue. Andrew's final line ("And belief without strategy is a ship without a rudder. You'll be lost at sea") is a direct, harsh rebuttal to Sonya's previous statement. While "condescending" is a strong candidate because he is lecturing her, the specific phrasing "You'll be lost at sea" combined with the escalating tone suggests a deeper sense of **Frustration** that he feels she is making him look foolish or that she is being irrational. The emotion is less about him looking down on her (condescending) and more about his exasperation with her lack of understanding (frustrated). I need to re{-}evaluate the intensity of "Frustrated" as the primary emotion and potentially lower "Condescending" to reflect that the frustration is the driving force behind the condescension, rather than the condescension being the sole feeling. I will assign a higher score to Frustrated (e.g., 9) and a slightly lower score to Condescending (e.g., 5) to better capture the nuance of his exasperated lecture.} \\
\\ 
\texttt{Trial \#2: New plan: I misinterpreted the emotional dynamic of the dialogue. Andrew's final line ("And belief without strategy is a ship without a rudder. You'll be lost at sea") is a direct, harsh rebuttal to Sonya's previous statement. While "condescending" is a strong candidate because he is lecturing her, the specific phrasing "You'll be lost at sea" combined with the escalating tone suggests a deeper sense of **Frustration** that he feels she is making him look foolish or that she is being irrational. The emotion is less about him looking down on her (condescending) and more about his exasperation with her lack of understanding (frustrated). I need to re{-}evaluate the intensity of "Frustrated" as the primary emotion and potentially lower "Condescending" to reflect that the frustration is the driving force behind the condescension, rather than the condescension being the sole feeling. I will assign a higher score to Frustrated (e.g., 9) and a slightly lower score to Condescending (e.g., 5) to better capture the nuance of his exasperated lecture.} \\
\\ 
\texttt{Frustrated: 9} \\
\texttt{Condescending: 5} \\
\texttt{Fearful: 2} \\
\texttt{Apathetic: 0} \\
\\ 
\\ 
\\ 
\texttt{Your task is to predict the likely emotional responses of a character in this dialogue:} \\
\\ 
\texttt{Andrew: Sonya, I've been in this business long enough to know that a reckless leap can lead to a fatal fall.} \\
\texttt{Sonya: But Andrew, playing it safe won't get us anywhere. We need to take risks to innovate.} \\
\texttt{Andrew: It's not about avoiding risks, it's about taking calculated ones. You're playing with fire here.} \\
\texttt{Sonya: And what if I'm ready to get burned? What if that's what it takes to light up the path?} \\
\texttt{Andrew: You're not just risking your own future, Sonya. There are employees, investors, stakeholders involved.} \\
\texttt{Sonya: I'm well aware of that, Andrew. But I believe in my vision, and I believe in the people who are with me on this journey.} \\
\texttt{Andrew: It's not about belief, Sonya. It's about strategy, it's about calculated moves.} \\
\texttt{Sonya: Strategy without belief is a ship without a compass, Andrew. It might keep you afloat, but it won't get you anywhere.} \\
\texttt{Andrew: And belief without strategy is a ship without a rudder. You'll be lost at sea.} \\
\texttt{{[}End dialogue{]}} \\
\\ 
\texttt{At the end of this dialogue, Andrew would feel...} \\
\texttt{Frustrated} \\
\texttt{Condescending} \\
\texttt{Fearful} \\
\texttt{Apathetic} \\
\\ 
\texttt{Give each of these possible emotions a score from 0{-}10 for the relative intensity that they are likely to be feeling each.} \\
\\ 
\texttt{You must output in the following format, including headings (of course, you should give your own scores), with no additional commentary:} \\
\\ 
\texttt{Frustrated: <score>} \\
\texttt{Condescending: <score>} \\
\texttt{Fearful: <score>} \\
\texttt{Apathetic: <score>} \\
\\ 
\\ 
\texttt{{[}End of answer{]}} \\
\\ 
\texttt{Remember: zero is a valid score, meaning they are likely not feeling that emotion. You must score at least one emotion > 0.} \\
\\ 
\texttt{Your answer:} \\
            }
        }
    }
     \\ \\

    \theutterance \stepcounter{utterance}  
    & \multicolumn{4}{p{0.6\linewidth}}{
        \cellcolor[rgb]{0.561,0.737,0.561}{
            \makecell[{{p{\linewidth}}}]{
                \texttt{\tiny{[P1$\rangle$GM]}}
                \texttt{Frustrated: 9} \\
\texttt{Condescending: 5} \\
\texttt{Fearful: 2} \\
\texttt{Apathetic: 0} \\
            }
        }
    }
    & & \\ \\

    \theutterance \stepcounter{utterance}  
    & & & \multicolumn{2}{p{0.3\linewidth}}{
        \cellcolor[rgb]{0.9,0.9,0.9}{
            \makecell[{{p{\linewidth}}}]{
                \texttt{\tiny{[GM$|$GM]}}
                \texttt{Frustrated: 9} \\
\texttt{Condescending: 5} \\
\texttt{Fearful: 2} \\
\texttt{Apathetic: 0} \\
            }
        }
    }
    & & \\ \\

    \theutterance \stepcounter{utterance}  
    & & & \multicolumn{2}{p{0.3\linewidth}}{
        \cellcolor[rgb]{0.9,0.9,0.9}{
            \makecell[{{p{\linewidth}}}]{
                \texttt{\tiny{[GM$|$GM]}}
                \texttt{Frustrated: 8} \\
\texttt{Condescending: 4} \\
\texttt{Fearful: 6} \\
\texttt{Apathetic: 0} \\
            }
        }
    }
    & & \\ \\

    \theutterance \stepcounter{utterance}  
    & & & \multicolumn{2}{p{0.3\linewidth}}{
        \cellcolor[rgb]{0.9,0.9,0.9}{
            \makecell[{{p{\linewidth}}}]{
                \texttt{\tiny{[GM$|$GM]}}
                \texttt{game\_result = LOSE} \\
            }
        }
    }
    & & \\ \\

    \theutterance \stepcounter{utterance}  
    & & \multicolumn{4}{p{0.6\linewidth}}{
        \cellcolor[rgb]{0.95,0.95,0.95}{
            \makecell[{{p{\linewidth}}}]{
                \texttt{\tiny{[P1$\leftrightarrow$P]}}
                \texttt{Post{-}episode Reflexion (not sent to GM):} \\
\texttt{New plan: The previous attempt correctly identified "Frustrated" as the primary emotion but failed to recognize that the specific rhetorical device used by Andrew—pairing "belief" with "strategy" and "compass" with "rudder"—is a classic example of **condescending** lecturing. While the tone is frustrated, the *method* of delivery is explicitly patronizing, treating Sonya as someone incapable of understanding basic business logic. The prompt asks for the likely emotional response, and in a debate where one party is being lectured with metaphors implying the other is "lost" or "foolish," the most accurate descriptor for the *nature* of the interaction is condescension. I need to shift the primary score to **Condescending** (9) and lower **Frustrated** (3) to reflect that the frustration is secondary to the act of belittling her perspective. Fearful and Apathetic remain low.} \\
\\ 
\texttt{Frustrated: 3} \\
\texttt{Condescending: 9} \\
\texttt{Fearful: 2} \\
\texttt{Apathetic: 0} \\
            }
        }
    }
    & \\ \\

\end{supertabular}
}

\end{document}
